\documentclass[12pt]{article}
\usepackage{float}
\usepackage{graphicx}
\usepackage{fontspec}
\usepackage{url}
\usepackage{apacite}
\setmainfont{Times New Roman}
\usepackage[letterpaper,top=2cm,bottom=2cm,left=2cm,right=2cm,marginparwidth=1.5cm]{geometry}

\title{Weekly assignment 10: Wrapup}
\author{Truong Khang Tobias Nguyen}
\date{}

\begin{document}

\maketitle

\section{Reflections}
As I come from a technical background and had no prior exposure to the non-technical part of cybersecurity such as organizational, legislative, or ethical aspects of cybersecurity, I would say that I did not have any expectations that the course would be technical. Because of that, I have learned a substantial amount of the organizational aspects of cybersecurity as well as standards, guidelines, legislations, and ethical principles that are relevant to cybersecurity.
\newline
The first part consisted of the standards and compliance part of cybersecurity. Although I am somewhat aware of these topics, e.g. I know that they exist, I never really put much thought into these topics and my assumption for both parts were "Experts develop standards and organizations implement them as a means to appear more trustworthy and an organization's employees should not comply to security measures for the sake of complying." Over the course of this part I have learned that these aspects are more nuanced than they may seem and that implementing standards is about considering its trade-offs and how well they fit an organization while compliance is a cultural question that questions on how organizations actually view cybersecurity.
\newline
Considering that I have already written two assignments about the two lectures about the legislative aspects of cybersecurity, I do not see the point of going too much into detail with it. What I can say is that I initially dreaded this aspects as I viewed legislation as an extremly boring topic that I did not really want to engage with. But Jens made the lectures quite fascinating which has been a pleasant surprise.
\newline
Finally comes the last part that concerned the ethical aspects. Even though I unfortunately missed the second and third lecture of this part, I do have to say that I easily enjoyed this part the most, especially the weekly assignments. Considering the technical focus of my Bachelor's degree, I never had any coursework on ethics in any capacity - in combination with my experience in software development that never required me to think about the ethical implications of what I develop (I suppose it would be different if I worked at a company like Palantir though), I never considered ethics in regards to my studies or work. This thought has been challenged with the first assignment of this part where we had to reflect on our own ethical stance, our values, and why it is ultimately important to consider ethics in our field. As these topics were never on my radar, it forced me to actually consider what I think is "good" or "bad" which has been an enjoyable experience because it made me actually reason on my own opinions instead of making points based on objective truths that are common in the technical fields. That being said, I found this part of the course the most interesting and the assignments for this part have also been the most enjoyable as they required me to solve questions based on my own stance which is also why I do believe that it is unfortunate that I have missed lectures of this part.
\newline
Regarding the part if this course has changed me as a cybersecurity student is difficult to say. On one hand, I could say "no" since I did not study cybersecurity up to this point and in this sense, this course has instead shaped my baseline views on cybersecurity rather than to change prior ones as I did not really have any beforehand. On the other hand, I could also say "yes" in the sense that this course changed the prior assumptions I had on cybersecurity as I come from a different background where cybersecurity was not covered. I do not believe that either point invalidates the other as they are not mutually exclusive and like most topics we have covered, I suppose that my answer is more nuanced. I would say that the lectures and assignments of this course (and the other courses obviously) have challenged my prior assumptions of cybersecurity and subsequently, took part of what has become the views on cybersecurity that I have now. 
\newline
To make an example, even though I have initially acknowledged the statement "Cybersecurity is not just a question of technical nature anymore but also a question of human nature.", I brushed the question off with the explaination that goes along the lines of "Cybersecurity obviously has to be a question of both technical and human nature, because even the most sophisticated technical solution can't protect an organization from human stupidity.". Even though I still stand by this statement, I acknowledge the fact that this statement is somewhat naive as it explains the need of the so-called "human firewall" because of human idiocy instead of asking why humans are stupid and how one can work around that. 
\newline 
To keep it short, the course has deepened my perspective on cybersecurity and why it actually needs to combine human and technical factors instead of treating humans as adversaries.

\section{Advice to Myself}
I am aware that I have entered the course with previous assumptions which do not sufficiently cover the nuances of organizational cybersecurity but I do not necessarily see it as a bad thing because I have learned quite a lot about because the content has challenged these assumptions which would not be the case if I started the course with that knowledge. If I entered the course with the knowledge I have now, I, quite frankly, would have grown bored of the course real quick and subsequently skipped most of the lectures as I would have found it unnecessary to attend them if I was already aware of it's contents.
\newline
That being said, there is not that much that I would tell my self of the past - instead, I would reinforce what I told myself at the start of the course which is "Be open to new ideas, let them challenge you, and learn from them". Just for fun, I would have told myself to not go into the legislative parts with dread, but to give this topic (albeit still a bit boring in my opinion) a chance at it is still important to know in the later parts of my career.

\section{Conclusion}
While I am still uncertain about my plans after finishing my studies, whether I actually want to pursue a PhD or if I want to continue my path as a software developer and although this course has been interesting and has exceeded my expectations, I am aware that I do not want to pursue a career in this field (no offense, really).
\newline
But for me, the most important thing about this course and it's effects on my later career is that I am aware of the topics covered by the course and it's implications and I will do my best to remember what I have learned here whenever I have to.

\end{document} 